\section{Conclusion}

In this paper, we widen the naming problem studies to the rule-based graph
transformation modeling systems. %
Using the formalism of both generalized maps and graph transformation rules, we
tackle the reevaluation mechanism task. %
While generalized maps homogeneously represent objects in all dimensions,
Jerboa's rules define geometric modeling operations we can easily syntactically
analyse. %
% We implement: a persistent name scheme where each persistent name represents a
% unique dart's history through the successive applications of rules and their
% matching nodes. %
First, we implement a persistent naming scheme. %
Each persistent name represents a unique dart's history through the successive
applications of rules and their matching nodes. %
% Then, for each persistent name, an evaluation DAG is built in order to trace an
% orbit's history from the bottom and up to the orbits it originates from. %
Then, for each persistent name, we build an evaluation DAG which traces an
orbit's history from the bottom and up to the orbits it originates from. %
To our knowledge, unlike other methods, our solution tracks only the entities
used in the parametric specification and the ones they originate from. %
Representing the complete history of an orbit in an evaluation DAG allows for an
efficient persistent naming mechanism that takes into account the impact of both
origins and traces modifications during reevaluation. %
Finally, reevaluation DAGs are built from a top-down traversal of evaluation
DAGs and allow matching each topological parameter on one or more, sometimes
none, values depending on the editing of the parametric specification. %
Thanks to our method, not only the naming problem is tackled within the usual
framework of parameters edition, but we also take the specification edition into
account (operation addition, deletion or move). Moreover, this approach provides
all the possible updated values of the parameters and, thus, enables
implementing different strategies. %

More complex operations can make use of several rules 
brought together
in a script. %
Later works will revolve around widening this reevaluation mechanism
to scripts. %

%%% Local Variables:
%%% mode: latex
%%% TeX-master: "../main"
%%% End:
